% ====================================================================
% FF3-4  fig:relaxode 재설계 — 기억 적분의 실수치 완화 (Ch1 문맥, §9.1 대체 제안)
%   좌표 = eq:lag 인과 기억 적분의 수치 평가 (gen_coords_ff3.py;
%   목표 ξ_eq(q)=logistic(중심 1.2, 폭 0.35), L_q=0.4 — 게이트: 0≤ξ_lag≤ξ_eq≤1, r≥0).
%   원 그림의 정보 요소(목표/지연 곡선·r 화살표·수직 점선·peak shape 주석) 전량 보존,
%   ξ축 0–1 실눈금·r(q) 곡선·기억 커널 음영·L_q 치수 추가.
%   사용 라이브러리: arrows.meta, calc, backgrounds (Ch1 preamble 내).
% ====================================================================
\begin{figure}[t]
\centering
\begin{tikzpicture}[x=3.2cm,y=4.6cm]
% ---- memory kernel shading at evaluation point q0=2.2 (background) ----
\begin{scope}[on background layer]
\fill[black!10] plot[smooth] coordinates {(0.6,0.004029) (0.68,0.004922) (0.76,0.006011) (0.84,0.007342) (0.92,0.008968) (1,0.01095) (1.08,0.01338) (1.16,0.01634) (1.24,0.01996) (1.32,0.02438) (1.4,0.02977) (1.48,0.03637) (1.56,0.04442) (1.64,0.05425) (1.72,0.06626) (1.8,0.08093) (1.88,0.09885) (1.96,0.1207) (2.04,0.1475) (2.12,0.1801) (2.2,0.22)} -- (2.2,0) -- (0.6,0) -- cycle;
\end{scope}
\draw[black!45] (2.2,0) -- (2.2,0.24);
\node[font=\tiny,anchor=west,align=left] at (2.28,0.175)
  {기억 커널 $e^{-(q-u)/L_q}$\\(평가점 $q{=}2.2$, $0.22$ 배율)};
\draw[{Bar[width=1.1mm]}-{Bar[width=1.1mm]}] (1.8,0.055) -- (2.2,0.055);
\node[font=\tiny,anchor=south] at (2.0,0.062) {$L_q$};
% ---- axes ----
\draw[-{Latex[length=1.4mm]}] (-0.05,0) -- (3.55,0)
  node[below,font=\scriptsize] {$q$ (capacity coord)};
\draw[-{Latex[length=1.4mm]}] (-0.05,0) -- (-0.05,1.13) node[above,font=\scriptsize] {$\xi$};
\foreach \yy/\lab in {0/0, 0.5/0.5, 1.0/1}
  {\draw (-0.03,\yy) -- (-0.07,\yy) node[left,font=\scriptsize] {\lab};}
\foreach \xx in {1,2,3} {\draw (\xx,0.008) -- (\xx,-0.008) node[below,font=\scriptsize] {\xx};}
\draw[densely dotted,black!50] (-0.05,1.0) -- (3.45,1.0);
% ---- target xi_eq (dashed) ----
\draw[densely dashed,thick] plot[smooth] coordinates {(0,0.03141) (0.1,0.04137) (0.2,0.05431) (0.3,0.071) (0.4,0.09231) (0.5,0.1192) (0.6,0.1526) (0.7,0.1933) (0.8,0.2418) (0.9,0.2979) (1,0.3609) (1.1,0.4291) (1.2,0.5) (1.3,0.5709) (1.4,0.6391) (1.5,0.7021) (1.6,0.7582) (1.7,0.8067) (1.8,0.8474) (1.9,0.8808) (2,0.9077) (2.1,0.929) (2.2,0.9457) (2.3,0.9586) (2.4,0.9686) (2.5,0.9762) (2.6,0.982) (2.7,0.9864) (2.8,0.9898) (2.9,0.9923) (3,0.9942) (3.1,0.9956) (3.2,0.9967) (3.3,0.9975) (3.4,0.9981)};
\node[font=\scriptsize,anchor=south] at (2.75,1.005) {target $\xi_\eq$};
% ---- lagged xi_lag (solid) : pointwise memory integral eq:lag ----
\draw[thick] plot[smooth] coordinates {(0,0.01482) (0.1,0.01959) (0.2,0.02584) (0.3,0.03399) (0.4,0.04454) (0.5,0.05809) (0.6,0.07533) (0.7,0.09698) (0.8,0.1237) (0.9,0.1562) (1,0.1947) (1.1,0.2392) (1.2,0.2894) (1.3,0.3442) (1.4,0.4023) (1.5,0.462) (1.6,0.5217) (1.7,0.5798) (1.8,0.6348) (1.9,0.6858) (2,0.7321) (2.1,0.7735) (2.2,0.8099) (2.3,0.8415) (2.4,0.8686) (2.5,0.8917) (2.6,0.9111) (2.7,0.9273) (2.8,0.9408) (2.9,0.9519) (3,0.9611) (3.1,0.9686) (3.2,0.9747) (3.3,0.9797) (3.4,0.9837)};
\node[font=\scriptsize,anchor=north west] at (1.98,0.70) {lagged $\xi_\mathrm{lag}$};
% ---- deficit r(q) = xi_eq - xi_lag (thin gray curve) ----
\draw[black!55] plot[smooth] coordinates {(0,0.01659) (0.1,0.02178) (0.2,0.02847) (0.3,0.03701) (0.4,0.04778) (0.5,0.06111) (0.6,0.07728) (0.7,0.09634) (0.8,0.1181) (0.9,0.1418) (1,0.1662) (1.1,0.1898) (1.2,0.2106) (1.3,0.2268) (1.4,0.2368) (1.5,0.24) (1.6,0.2365) (1.7,0.2269) (1.8,0.2126) (1.9,0.195) (2,0.1755) (2.1,0.1555) (2.2,0.1358) (2.3,0.1171) (2.4,0.09997) (2.5,0.08456) (2.6,0.07096) (2.7,0.05914) (2.8,0.04899) (2.9,0.04037) (3,0.03311) (3.1,0.02705) (3.2,0.02201) (3.3,0.01785) (3.4,0.01444)};
\fill[black!55] (1.5,0.24) circle (0.9pt);
\node[black!55,font=\tiny,anchor=south west,align=left] at (0.10,0.245)
  {$r(q)=\xi_\eq-\xi_\mathrm{lag}$ (최대 $0.240$ @ $q{=}1.5$)};
% ---- gap arrow at q=1.5 ----
\draw[{Latex[length=1.2mm]}-{Latex[length=1.2mm]},thick] (1.5,0.462) -- (1.5,0.7021);
\node[font=\scriptsize,anchor=west] at (1.54,0.545) {$r=\xi_\eq-\xi_\mathrm{lag}$};
\draw[densely dotted,gray] (1.5,0) -- (1.5,0.462);
% ---- peak-shape annotation ----
\node[font=\scriptsize,anchor=west] at (0.05,0.93)
  {peak shape $=(\xi_\eq-\xi_\mathrm{lag})/L_V$};
\end{tikzpicture}
\caption{지연의 완화 --- 좌표는 식~\eqref{eq:lag} 기억 적분의 수치 평가다(목표
$\xi_\eq(q)$ 는 중심 $1.2$$\cdot$폭 $0.35$ 의 logistic, $L_q{=}0.4$). 파선 $=$ 평형 목표 $\xi_\eq$,
굵은 실선 $=$ 그 과거를 규격화 지수 커널로 가중평균한 지연 진행률 $\xi_\mathrm{lag}$ --- 어느 $q$
에서나 $0\le\xi_\mathrm{lag}\le\xi_\eq\le1$ 이다. 회색 곡선이 결핍 $r(q)=\xi_\eq-\xi_\mathrm{lag}$
로, 목표가 가장 빨리 변하는 중심 부근에서 최대($0.240$ @ $q{=}1.5$, 화살표)가 되고 peak 모양
$(\xi_\eq-\xi_\mathrm{lag})/L_V$(식~\eqref{eq:peakshape})는 이 차를 길이로 나눈 것이다. 오른쪽 음영이
평가점 $q{=}2.2$ 의 기억 커널 $e^{-(q-u)/L_q}$(표시 배율 $0.22$) --- 계는 길이 $L_q$ 규모의 최근
과거만 기억하며, $L_{V}\to0$ 에서 이 차분이 평형 종의 미분으로 수렴하는 극한은
식~\eqref{eq:tail-limit} 이 닫는다.}
\label{fig:cand-ff3-4}
\end{figure}
